\subsection{迭代导数}

1.7.1. 设 \(f\) 在 E 中一点 \(x_{0}\) 的某邻域上定义且取值于 F。若 \(f\) 在 \(x_{0}\) 的某邻域上可微，则其导数 Df 是从 \(x_{0}\) 的某邻域到多范数空间 \(\mathcal{L}(\mathrm{E};\mathrm{F})\) 的映射，其中后者是从 E 到 F 的连续线性映射空间。设 \(p\) 是满足 \(\geqslant 2\) 的整数：

若 \(f\)\(p\) 在 \(x_0\) 的某邻域上可微，且其导数 \(f\) 在 \(x_0\) 处 \(Df\) 次可微，则称 \((p - 1)\) 在 \(x_0\) 处 \(p\) 次可微。于是定义 \(f\) 在 \(x_0\) 处的 \(p\) 阶导数；这是由下式定义的从 \(D^p f(x_0)\) 到 F 的连续 \(E^p\) 线性映射 \(F\)：

\[
\mathrm{D} ^ {p} f (x _ {0}). (h _ {1}, \dots , h _ {p}) = (\mathrm{D} (\mathrm{D} ^ {p - 1} f) (x _ {0}). h _ {1}). (h _ {2}, \dots , h _ {p}).
\]

此外置 \(D^{0}f = f\) 且 \(D^{1}f = Df\)。若 \(f\) 在 \(p\) 处 \(x_{0}\) 次可微，且 \(q\)、\(s\) 是满足 \(q + s = p\) 且 \(s > 0\) 的两个整数，则 \(f\) 在 \(q\) 的某邻域上 \(x_{0}\) 次可微，取值于 \(D^{q}f\) 的函数 \(\mathcal{L}_{q}(E; F)\) 在 \(s\) 处 \(x_{0}\) 次可微，并且有：

\[
\mathrm{D} ^ {q + s} f (x _ {0}). (h _ {1}, \dots , h _ {q + s}) = (\mathrm{D} ^ {s} (\mathrm{D} ^ {q} f) (x _ {0}). (h _ {1}, \dots , h _ {s})) \cdot (h _ {s + 1}, \dots , h _ {q + s})
\]

这一关系按记号滥用写成：

\[
\mathbf {D} ^ {q + s} f = \mathbf {D} ^ {s} \mathbf {D} ^ {q} f.
\]

1.7.2. 设 \((\mathbf{E}_{i})_{1\leqslant i\leqslant n}\) 是 E 的闭向量子空间，且 E 是各 \(E_{i}\) 的拓扑直和。于是，在存在时定义从 \(D_{i_{1}}\ldots D_{i_{m}}f\) 的某邻域到 F 的映射 f 的迭代偏导数 \(x_{0}\in E\)；它是从

\[
\mathbf {E} _ {i _ {1}} \times \dots \times \mathbf {E} _ {i _ {m}}
\]

到 F 的连续多线性映射，按整数 \(m \geqslant 1\) 归纳定义如下：若 \(D_{i_{2}} \ldots D_{i_{m}} f(x)\) 在 \(x_{0}\) 的某邻域上存在，且关于 \(E_{i_{1}}\) 有偏导数，则 \(D_{i_{1}} \ldots D_{i_{m}} f(x_{0})\) 由下式给出：

\[
\mathrm{D} _ {i _ {1}} \dots \mathrm{D} _ {i _ {m}} f (x _ {0}). (h _ {1}, \dots , h _ {m}) = (\mathrm{D} _ {i _ {1}} (\mathrm{D} _ {i _ {2}} \dots \mathrm{D} _ {i _ {m}} f) (x _ {0}). h _ {1}). (h _ {2}, \dots , h _ {m})
\]

其中 \(h_{k} \in E_{i_{k}}\)。

若 \(f\) 在 \(m\) 处 \(x_{0}\) 次可微，则偏导数 \(\mathbf{D}_{i_{1}}\ldots\mathbf{D}_{i_{m}}f(x_{0})\) 存在，并等于 \(\mathbf{D}^{m}f(x_{0})\) 对 \(\mathbf{E}_{i_{1}}\times\cdots\times\mathbf{E}_{i_{m}}\) 的子空间 \(\mathbf{E}^{m}\) 的限制。因此，\(\mathbf{D}^{m}f(x_{0})\) 完全由 \(m\) 处 \(x_{0}\) 阶迭代偏导数确定。

1.7.3. 假设 E 是有限维的，且 \((e_{1},\ldots,e_{n})\) 是 E 的一组基。置 \(E_{i}=Ke_{i}\)，并设 f 是从 \(x_{0}\) 的某邻域到 F 的映射。若偏导数 \(D_{i_{1}}\ldots D_{i_{m}}f(x_{0})\)（使用 1.7.2 的记号）存在，则置：

\[
\partial_ {i _ {1}} \dots \partial_ {i _ {m}} f (x _ {0}) = \mathbf {D} _ {i _ {1}} \dots \mathbf {D} _ {i _ {m}} f (x _ {0}). (e _ {i _ {1}}, \dots , e _ {i _ {m}}).
\]

